\section{The arithmetic interface}
\label{sec:arithmetic}

The preceding results separate information about a reflected
divisor from information about an arithmetic formula. For
clarity, we end by writing that interface in exactly the
normalization used in this article.

Let \(Z_\zeta=\{\rho-\tfrac12\}\), where \(\rho\) runs over
the nontrivial zeros of the Riemann zeta function, with
multiplicity. The functional equation supplies reflection,
and the classical ordinate count is polynomial. For
\(f\in C_c^\infty(-L,L)\), write
\[
 F_f(t)=H_f(it)=\int_\R f(u)e^{itu}\,du,\qquad
 C_f(v)=\int_\R f(u)\overline{f(u-v)}\,du.
\]
Thus \(C_f(-v)=\overline{C_f(v)}\) and
\(C_f(0)=\|f\|_2^2\). The Guinand--Weil explicit formula,
in additive coordinates, is
\begin{equation}\label{arith:actual-form}
 \begin{split}
 Q_{Z_\zeta}(f,f)
 ={}&\frac1{2\pi}\int_\R
  \left\{\Re\psi\left(\frac14+\frac{it}2\right)-\log\pi\right\}
                     |F_f(t)|^2\,dt\\
 &+2\Re\left\{H_f(1/2)\overline{H_f(-1/2)}\right\}\\
 &-2\Re\sum_{n\ge2}\frac{\Lambda(n)}{\sqrt n}\,C_f(\log n).
 \end{split}
\end{equation}
Here \(\psi=\Gamma'/\Gamma\), and \(\Lambda(n)\) equals
\(\log p\) if \(n\) is a positive power of a prime \(p\)
and equals zero otherwise. See
\cite[Sections 2--4]{Bombieri} for the explicit formula
and its associated quadratic functional. Our formula follows
by applying it to \(f*\widetilde f\), where
\(\widetilde f(u)=\overline{f(-u)}\): its transform is
\(H_f(z)\overline{H_f(-\overline z)}\), exactly the
product in \eqref{eq:form}. The Fourier convention accounts
for the factor \(1/(2\pi)\).

Every term in \eqref{arith:actual-form} is well defined on
the stated class. The prime sum is finite because
\(C_f(v)=0\) for \(|v|\ge2L\). The Fourier transform decreases
rapidly, while the real gamma multiplier grows like
\(\log|t|\) and is bounded below. The second line is a
polar contribution, not a nonnegative form. In particular,
for a real odd \(f\) it equals \(-2|H_f(1/2)|^2\).
Deleting that line or replacing it by a positive quantity
would change the arithmetic problem.

Here is the exact implication supplied by the abstract
lower-growth result. Suppose that, for one fixed
\(0\le\kappa<1\) and every \(\varepsilon>0\), an independent
arithmetic argument proves
\begin{equation}\label{arith:open-bound}
 \cM_{Z_\zeta}(L)\le
 C_\varepsilon e^{(\kappa+\varepsilon)L}
 \qquad(L\ge1).
\end{equation}
If there is an off-axis zero, the lower rate in
Section~\ref{sec:growth} gives
\[
 2\sup_\rho\left|\Re\rho-\tfrac12\right|
 \le\liminf_{L\to\infty}\frac{\log\cM_{Z_\zeta}(L)}L
 \le\kappa.
\]
If there is no off-axis zero, the same zero-location
conclusion is immediate. Hence
\begin{equation}\label{arith:strip}
 \left|\Re\rho-\tfrac12\right|\le\kappa/2
 \quad\hbox{for every nontrivial zero } \rho.
\end{equation}
In particular, subexponential negative growth would imply
the Riemann Hypothesis. The weighted sufficient condition
of Section~\ref{sec:growth} is not needed for this
implication; only the count-based lower rate is used.

The hypothesis \eqref{arith:open-bound} is not proved in
this article. The signed prime sum, the gamma term, and
the polar contribution would have to be controlled
together, uniformly over every admissible test in the
window. Counting bounds for a divisor and positivity of
one finite compression do not provide that estimate.
Conversely, the arbitrary-growth divisors of
Section~\ref{sec:arbitrary-growth} are not asserted to be
zero sets of arithmetic \(L\)-functions and do not refute
an arithmetic inequality.

Several questions remain on the abstract side as well.
The constants in the support--margin law are not uniform
in a growing count budget; the exact leading minimax
constant is undetermined. The weighted local bound is a
sufficient condition for the matching upper rate; it is
not asserted to be necessary. A characterization of growth that
retains the compensation supplied by the axis part
would be stronger than either a pure counting bound
or an absolute estimate of the off-axis contribution.
These are distinct questions from proving
\eqref{arith:open-bound}.
